\documentclass[12pt,letterpaper]{article}
\usepackage{graphicx,textcomp}
\usepackage{natbib}
\usepackage{setspace}
\usepackage{fullpage}
\usepackage{color}
\usepackage[reqno]{amsmath}
\usepackage{amsthm}
\usepackage{fancyvrb}
\usepackage{amssymb,enumerate}
\usepackage[all]{xy}
\usepackage{endnotes}
\usepackage{lscape}
\newtheorem{com}{Comment}
\usepackage{float}
\usepackage{hyperref}
\newtheorem{lem} {Lemma}
\newtheorem{prop}{Proposition}
\newtheorem{thm}{Theorem}
\newtheorem{defn}{Definition}
\newtheorem{cor}{Corollary}
\newtheorem{obs}{Observation}
\usepackage[compact]{titlesec}
\usepackage{dcolumn}
\usepackage{tikz}
\usetikzlibrary{arrows}
\usepackage{multirow}
\usepackage{xcolor}
\newcolumntype{.}{D{.}{.}{-1}}
\newcolumntype{d}[1]{D{.}{.}{#1}}
\definecolor{light-gray}{gray}{0.65}
\usepackage{url}
\usepackage{listings}
\usepackage{color}

\definecolor{codegreen}{rgb}{0,0.6,0}
\definecolor{codegray}{rgb}{0.5,0.5,0.5}
\definecolor{codepurple}{rgb}{0.58,0,0.82}
\definecolor{backcolour}{rgb}{0.95,0.95,0.92}

\lstdefinestyle{mystyle}{
	backgroundcolor=\color{backcolour},   
	commentstyle=\color{codegreen},
	keywordstyle=\color{magenta},
	numberstyle=\tiny\color{codegray},
	stringstyle=\color{codepurple},
	basicstyle=\footnotesize,
	breakatwhitespace=false,         
	breaklines=true,                 
	captionpos=b,                    
	keepspaces=true,                 
	numbers=left,                    
	numbersep=5pt,                  
	showspaces=false,                
	showstringspaces=false,
	showtabs=false,                  
	tabsize=2
}
\lstset{style=mystyle}
\newcommand{\Sref}[1]{Section~\ref{#1}}
\newtheorem{hyp}{Hypothesis}

\title{Problem Set 2}
\date{Due: February 18, 2026}
\author{Applied Stats II}


\begin{document}
	\maketitle
	\section*{Instructions}
	\begin{itemize}
		\item Please show your work! You may lose points by simply writing in the answer. If the problem requires you to execute commands in \texttt{R}, please include the code you used to get your answers. Please also include the \texttt{.R} file that contains your code. If you are not sure if work needs to be shown for a particular problem, please ask.
		\item Your homework should be submitted electronically on GitHub in \texttt{.pdf} form.
		\item This problem set is due before 23:59 on Wednesday February 18, 2026. No late assignments will be accepted.

	\end{itemize}

	\vspace{.25cm}

\noindent We're interested in what types of international environmental agreements or policies people support (\href{https://www.pnas.org/content/110/34/13763}{Bechtel and Scheve 2013)}. So, we asked 8,500 individuals whether they support a given policy, and for each participant, we vary the (1) number of countries that participate in the international agreement and (2) sanctions for not following the agreement. \\

\noindent Load in the data labeled \texttt{climateSupport.RData} on GitHub, which contains an observational study of 8,500 observations.

\begin{itemize}
	\item
	Response variable: 
	\begin{itemize}
		\item \texttt{choice}: 1 if the individual agreed with the policy; 0 if the individual did not support the policy
	\end{itemize}
	\item
	Explanatory variables: 
	\begin{itemize}
		\item
		\texttt{countries}: Number of participating countries [20 of 192; 80 of 192; 160 of 192]
		\item
		\texttt{sanctions}: Sanctions for missing emission reduction targets [None, 5\%, 15\%, and 20\% of the monthly household costs given 2\% GDP growth]
		
	\end{itemize}
	
\end{itemize}

\newpage
\noindent Please answer the following questions:

\begin{enumerate}
	\item
	Remember, we are interested in predicting the likelihood of an individual supporting a policy based on the number of countries participating and the possible sanctions for non-compliance.
	\begin{enumerate}
		\item [] Fit an additive model. Provide the summary output, the global null hypothesis, and $p$-value. Please describe the results and provide a conclusion.

	\end{enumerate}
	
	\item
	If any of the explanatory variables are statistically significant in this model, then:
	\begin{enumerate}
		\item
		For the policy in which nearly all countries participate [160 of 192], how does increasing sanctions from 5\% to 15\% change the odds that an individual will support the policy? (Interpretation of a coefficient)
		\item
		For the policy in which very few countries participate [20 of 192], how does increasing sanctions from 5\% to 15\% change the odds that an individual will support the policy? (Interpretation of a coefficient)
		\item
		What is the estimated probability that an individual will support a policy if there are 80 of 192 countries participating with no sanctions? 
	
	\end{enumerate}
	\item
	Would the answers to 2a and 2b potentially change if we included an interaction term in this model? Why? 
	\begin{itemize}
		\item Perform a test to see if including an interaction is appropriate.
	\end{itemize}
	\end{enumerate}

\newpage

\section*{Problem 1}

\noindent\textbf{Model and reference categories.}
I estimated a logistic regression of support for the policy on \texttt{countries} and \texttt{sanctions}, treating both explanatory variables as \emph{unordered categorical} factors. The reference categories are \texttt{countries = ``20 of 192''} and \texttt{sanctions = ``None''}. Therefore, the intercept corresponds to the log-odds of support for the baseline scenario (20 of 192, None), and each reported coefficient is interpreted relative to that baseline holding the other variable constant.

\vspace{0.75em}
\noindent\textbf{Global null hypothesis test.}
The global null is that all slope coefficients are zero (i.e., neither \texttt{countries} nor \texttt{sanctions} is associated with support). Using a likelihood-ratio (deviance) test, I obtain $\chi^2(5)=215.15$ with $p=1.62\times 10^{-44}$. I reject the global null; at least one category differs from the reference in support.

\vspace{1.0em}
\noindent\textbf{Regression output (additive model).}
% --- Paste stargazer output for model1 below (generated in R) ---
\begin{table}[!htbp] \centering 
	\small
	\caption{Logistic Regression: Support for International Environmental Policy} 
	\label{tab:model1} 
	\begin{tabular}{@{\extracolsep{5pt}}lc} 
		\\[-1.8ex]\hline 
		\hline \\[-1.8ex] 
		& \multicolumn{1}{c}{\textit{Dependent variable:}} \\ 
		\cline{2-2} 
		\\[-1.8ex] & choice \\ 
		\hline \\[-1.8ex] 
		countries80 of 192 & 0.336$^{***}$ \\ 
		& (0.054) \\ 
		& \\ 
		countries160 of 192 & 0.648$^{***}$ \\ 
		& (0.054) \\ 
		& \\ 
		sanctions5\% & 0.192$^{***}$ \\ 
		& (0.062) \\ 
		& \\ 
		sanctions15\% & $-$0.133$^{**}$ \\ 
		& (0.062) \\ 
		& \\ 
		sanctions20\% & $-$0.304$^{***}$ \\ 
		& (0.062) \\ 
		& \\ 
		Constant & $-$0.273$^{***}$ \\ 
		& (0.054) \\ 
		& \\ 
		\hline \\[-1.8ex] 
		Observations & 8,500 \\ 
		Log Likelihood & $-$5,784.130 \\ 
		Akaike Inf. Crit. & 11,580.260 \\ 
		\hline 
		\hline \\[-1.8ex] 
		\textit{Note:}  & \multicolumn{1}{r}{$^{*}$p$<$0.1; $^{**}$p$<$0.05; $^{***}$p$<$0.01} \\ 
	\end{tabular} 
\end{table} 

\newpage

\noindent\textbf{R code used:}

\lstinputlisting[language=R, firstline=42, lastline=85]{PS02_JL.R}



\section*{Problem 2}

\noindent Because there is \emph{no interaction term}, the effect of changing \texttt{sanctions} is the same regardless of the value of \texttt{countries}. Formally, for two sanctions levels $a$ and $b$:
\[
\log\!\left(\frac{\text{odds}(\text{Supported}\mid a)}{\text{odds}(\text{Supported}\mid b)}\right)
= \delta_{a}-\delta_{b},
\]
which does \textit{not} depend on \texttt{countries}.

\begin{enumerate}
	\item[(a)] \textbf{Countries = 160 of 192; sanctions 5\% $\to$ 15\%.} The change in log-odds is
	\[
	\Delta = \delta_{15}-\delta_{5} = -0.13325 - 0.19186 = -0.32510.
	\]
	The corresponding odds ratio is $e^{\Delta} \approx 0.722$. Thus, increasing sanctions from 5\% to 15\% multiplies the odds of support by about $0.72$ (a \(\approx 27.8\%\) decrease in odds), holding \texttt{countries} fixed at 160 of 192.
	
	\item[(b)] \textbf{Countries = 20 of 192; sanctions 5\% $\to$ 15\%.} In an additive model, this odds ratio is identical to part (a):
	\[
	OR = e^{\delta_{15}-\delta_{5}} \approx 0.722.
	\]
	Interpretation: increasing sanctions from 5\% to 15\% decreases the odds of support by about 27.8\%, holding \texttt{countries} fixed at 20 of 192.
	
	\item[(c)] \textbf{Predicted probability for 80 of 192 countries and no sanctions.} For \texttt{countries = 80 of 192} and \texttt{sanctions = None} (reference), the linear predictor is:
	\[
	\hat{\eta} = \hat{\beta}_0 + \hat{\beta}_{80} = -0.27266 + 0.33636 = 0.06370.
	\]
	Applying the logistic transform gives:
	\[
	\hat{p} = \frac{1}{1+e^{-\hat{\eta}}} \approx 0.5159.
	\]
	Therefore, the estimated probability of support is approximately \(\hat{p}\approx 0.516\) (51.6\%).
\end{enumerate}

\noindent\textbf{R code used:}

\lstinputlisting[language=R, firstline=98, lastline=138]{PS02_JL.R}


\newpage
\section*{Problem 3}

\begin{enumerate}
	\item[(i)] \textbf{Would the answers to 2(a) and 2(b) differ if we included an interaction?}\\
	Yes, potentially. With a \texttt{countries $\times$ sanctions} interaction, the effect of increasing sanctions could vary across the \texttt{countries} categories, so the odds ratio for 5\% $\to$ 15\% could differ between ``160 of 192'' and ``20 of 192''.
	
	\item[(ii)] \textbf{Test whether the interaction is appropriate}\\
	I compared the additive model to a model including \texttt{countries:sanctions} using a likelihood-ratio test. The deviance improvement is $6.2928$ on $6$ degrees of freedom ($p=0.3912$). I therefore do \emph{not} find evidence that the interaction improves model fit; the additive model is adequate, which supports why 2(a) and 2(b) are the same in my results.

\noindent\textbf{R code used:}

\lstinputlisting[language=R, firstline=145, lastline=151]{PS02_JL.R}
	
	% --- Paste stargazer output for model1 vs interaction below (generated in R) ---
	\begin{table}[!htbp] \centering
		\caption{Model Comparison: Additive vs. Interaction Model}
		\label{tab:model_comparison}
		\small
		\begin{tabular}{@{\extracolsep{5pt}}lcc}
			\\[-1.8ex]\hline
			\hline \\[-1.8ex]
			& \multicolumn{2}{c}{\textit{Dependent variable:}} \\
			\cline{2-3}
			\\[-1.8ex] & \multicolumn{2}{c}{choice} \\
			\\[-1.8ex] & (1) & (2)\\
			\hline \\[-1.8ex]
			countries80 of 192 & 0.336$^{***}$ & 0.376$^{***}$ \\
			& (0.054) & (0.106) \\
			& & \\
			countries160 of 192 & 0.648$^{***}$ & 0.613$^{***}$ \\
			& (0.054) & (0.108) \\
			& & \\
			sanctions5\% & 0.192$^{***}$ & 0.122 \\
			& (0.062) & (0.105) \\
			& & \\
			sanctions15\% & $-$0.133$^{**}$ & $-$0.097 \\
			& (0.062) & (0.108) \\
			& & \\
			sanctions20\% & $-$0.304$^{***}$ & $-$0.253$^{**}$ \\
			& (0.062) & (0.108) \\
			& & \\
			countries80 of 192:sanctions5\% &  & 0.095 \\
			&  & (0.152) \\
			& & \\
			countries160 of 192:sanctions5\% &  & 0.130 \\
			&  & (0.151) \\
			& & \\
			countries80 of 192:sanctions15\% &  & $-$0.052 \\
			&  & (0.152) \\
			& & \\
			countries160 of 192:sanctions15\% &  & $-$0.052 \\
			&  & (0.153) \\
			& & \\
			countries80 of 192:sanctions20\% &  & $-$0.197 \\
			&  & (0.151) \\
			& & \\
			countries160 of 192:sanctions20\% &  & 0.057 \\
			&  & (0.154) \\
			& & \\
			Constant & $-$0.273$^{***}$ & $-$0.275$^{***}$ \\
			& (0.054) & (0.075) \\
			& & \\
			\hline \\[-1.8ex]
			Observations & 8,500 & 8,500 \\
			Log Likelihood & $-$5,784.130 & $-$5,780.983 \\
			Akaike Inf. Crit. & 11,580.260 & 11,585.970 \\
			\hline
			\hline \\[-1.8ex]
			\textit{Note:} & \multicolumn{2}{r}{$^{*}$p$<$0.1; $^{**}$p$<$0.05; $^{***}$p$<$0.01} \\
		\end{tabular}
	\end{table}
\end{enumerate}


\end{document}
